% GENERATED FILE. Edit canonical lessons, then run npm run generate:books.
% canonical-source-hash: fnv1a64:eb231d3587b49682
% canonical-lessons: ZH-C299-dejiu, ZH-C299-jishi, ZH-C299-jijiu, ZH-C299-yundao, ZH-C299-qiangjiu

\chapter{Help Arrives}
\label{ch:zh-help-arrives}
\hlchaptermodality{299}

\begin{chapteropening}
\textbf{By the end of this chapter:} \emph{I can say people were rescued in time, talk about first aid, and say the doctors fought to save a life.}

Complete the last lesson of chapter 299: I can say people were rescued in time, talk about first aid, and say the doctors fought to save a life.
\end{chapteropening}

\section[\texorpdfstring{\zh{得救} (déjiù)}{得救 (déjiù)}]{\zh{得救} (déjiù) — to be rescued; to be saved from danger}

\label{lesson:ZH-C299-dejiu}



\begin{quote}
Before the new one: say the Mandarin for to shout loudly, then the Mandarin for the cry for help.
\end{quote}

\subsection*{You'll want to know: \zh{得救}}
\textbf{\zh{得救}} — \emph{déjiù} — \textquotedblleft{}to be rescued; to be saved from danger\textquotedblright{}.

\begin{quote}
\textbf{\zh{事故发生以后大家都得救了}} \emph{shìgù fāshēng yǐhòu, dàjiā dōu déjiù le} — After the accident, everyone was rescued.
\end{quote}

\begin{cousinweb}
\textbf{\zh{得救}} (\emph{déjiù}) is built from \zh{得} 'to get' and \zh{救} 'to save', both inherited from Middle Chinese, \emph{tok} and \emph{kjuwH} (Guangyun, in Baxter's transcription), which glosses the first 'to gain': to get saved.
\end{cousinweb}

\begin{grammarlens}[title={What you've built}]
One more word for telling what happened.
\end{grammarlens}

\subsection*{Your turn}
\begin{itemize}
\raggedright
  \item \emph{Say it:} \emph{déjiù}
  \item \emph{Say it:} \emph{déjiù}, once more
  \item \emph{Say it:} \emph{shìgù fāshēng yǐhòu, dàjiā dōu déjiù le} — After the accident, everyone was rescued.
  \item \emph{Recall:} say the Mandarin for to shout loudly, then the Mandarin for the cry for help, then say \emph{déjiù} again
\end{itemize}

\begin{tcolorbox}[breakable,colback=teal!4,colframe=teal!35!black,title={Before you move on}]
What does \zh{得救} mean? (\textquotedblleft{}To be rescued; to be saved from danger\textquotedblright{}.) Say it once more.
\end{tcolorbox}
\section[jíshí]{jíshí — in time; promptly; before it was too late}

\label{lesson:ZH-C299-jishi}



\begin{quote}
Before the new one: say the Mandarin for the cry for help, then the Mandarin for to be saved from danger.
\end{quote}

\subsection*{You'll want to know: jíshí}
\emph{jíshí} — \textquotedblleft{}in time; promptly; before it was too late\textquotedblright{}.

\begin{quote}
\emph{yīshēng lái de hěn jíshí} — The doctor came in good time.
\end{quote}

Its characters come later; for now it is a word for the ear.

\begin{cousinweb}
\emph{jíshí} is built from \emph{jí} 'to reach' and \zh{时} 'time', both inherited from Middle Chinese, \emph{gip} and \emph{dzyi} (Guangyun, in Baxter's transcription), which glosses the first 'to arrive' and 'to catch up with': to catch the time.
\end{cousinweb}

\begin{grammarlens}[title={What you've built}]
One more word for telling what happened.
\end{grammarlens}

\subsection*{Your turn}
\begin{itemize}
\raggedright
  \item \emph{Say it:} \emph{jíshí}
  \item \emph{Say it:} \emph{jíshí}, once more
  \item \emph{Say it:} \emph{yīshēng lái de hěn jíshí} — The doctor came in good time.
  \item \emph{Recall:} say the Mandarin for the cry for help, then the Mandarin for to be saved from danger, then say \emph{jíshí} again
\end{itemize}

\begin{tcolorbox}[breakable,colback=teal!4,colframe=teal!35!black,title={Before you move on}]
What does jíshí mean? (\textquotedblleft{}In time; promptly; before it was too late\textquotedblright{}.) Say it once more.
\end{tcolorbox}
\section[jíjiù]{jíjiù — first aid}

\label{lesson:ZH-C299-jijiu}



\begin{quote}
Before the new one: say the Mandarin for to be saved from danger, then the Mandarin for before it was too late.
\end{quote}

\subsection*{You'll want to know: jíjiù}
\emph{jíjiù} — \textquotedblleft{}first aid\textquotedblright{}.

\begin{quote}
\emph{wǒ xué guo jíjiù} — I've done a first-aid course.
\end{quote}

Its characters come later; for now it is a word for the ear.

\begin{cousinweb}
\emph{jíjiù} is built from \emph{jí} 'urgent' and \zh{救} 'to save', both inherited from Middle Chinese, \emph{kip} and \emph{kjuwH} (Guangyun, in Baxter's transcription), which glosses the first 'hurried'.
\end{cousinweb}

\begin{grammarlens}[title={What you've built}]
One more word for telling what happened.
\end{grammarlens}

\subsection*{Your turn}
\begin{itemize}
\raggedright
  \item \emph{Say it:} \emph{jíjiù}
  \item \emph{Say it:} \emph{jíjiù}, once more
  \item \emph{Say it:} \emph{wǒ xué guo jíjiù} — I've done a first-aid course.
  \item \emph{Recall:} say the Mandarin for to be saved from danger, then the Mandarin for before it was too late, then say \emph{jíjiù} again
\end{itemize}

\begin{tcolorbox}[breakable,colback=teal!4,colframe=teal!35!black,title={Before you move on}]
What does jíjiù mean? (\textquotedblleft{}First aid\textquotedblright{}.) Say it once more.
\end{tcolorbox}
\section[yūndǎo]{yūndǎo — to faint; to pass out}

\label{lesson:ZH-C299-yundao}



\begin{quote}
Before the new one: say the Mandarin for before it was too late, then the Mandarin for first aid.
\end{quote}

\subsection*{You'll want to know: yūndǎo}
\emph{yūndǎo} — \textquotedblleft{}to faint; to pass out\textquotedblright{}.

\begin{quote}
\emph{tā tūrán yūndǎo le} — She suddenly fainted.
\end{quote}

Its characters come later; for now it is a word for the ear.

\begin{cousinweb}
\emph{yūndǎo} is built from \emph{yūn} 'dizzy' and \emph{dǎo} 'to fall over', both inherited from Middle Chinese, \emph{hjunH} and \emph{tawX} (Guangyun, in Baxter's transcription), which glosses the first 'a halo round the sun or moon' and the second 'to fall flat'. The sense 'dizzy' is not given there, and the level tone of yūn is irregular for \emph{hjunH}.
\end{cousinweb}

\begin{grammarlens}[title={What you've built}]
One more word for telling what happened.
\end{grammarlens}

\subsection*{Your turn}
\begin{itemize}
\raggedright
  \item \emph{Say it:} \emph{yūndǎo}
  \item \emph{Say it:} \emph{yūndǎo}, once more
  \item \emph{Say it:} \emph{tā tūrán yūndǎo le} — She suddenly fainted.
  \item \emph{Recall:} say the Mandarin for before it was too late, then the Mandarin for first aid, then say \emph{yūndǎo} again
\end{itemize}

\begin{tcolorbox}[breakable,colback=teal!4,colframe=teal!35!black,title={Before you move on}]
What does yūndǎo mean? (\textquotedblleft{}To faint; to pass out\textquotedblright{}.) Say it once more.
\end{tcolorbox}
\section[qiǎngjiù]{qiǎngjiù — to fight to save a patient; to give emergency treatment}

\label{lesson:ZH-C299-qiangjiu}



\begin{quote}
Before the new one: say the Mandarin for first aid, then the Mandarin for to pass out.
\end{quote}

\subsection*{You'll want to know: qiǎngjiù}
\emph{qiǎngjiù} — \textquotedblleft{}to fight to save a patient; to give emergency treatment\textquotedblright{}.

\begin{quote}
\emph{yīshēng zhèngzài qiǎngjiù tā} — The doctors are fighting to save him.
\end{quote}

Its characters come later; for now it is a word for the ear.

\begin{cousinweb}
\emph{qiǎngjiù} is built from \emph{qiǎng} 'to snatch' and \zh{救} 'to save', both inherited from Middle Chinese, \emph{tshjangX} and \emph{kjuwH} (Guangyun, in Baxter's transcription). The Guangyun gives the first no gloss in this reading and glosses it 'to resist' and 'to dash' in the level tone; the sense 'to snatch' is not given there.
\end{cousinweb}

\begin{grammarlens}[title={What you've built}]
One more word for telling what happened.
\end{grammarlens}

\subsection*{Your turn}
\begin{itemize}
\raggedright
  \item \emph{Say it:} \emph{qiǎngjiù}
  \item \emph{Say it:} \emph{qiǎngjiù}, once more
  \item \emph{Say it:} \emph{yīshēng zhèngzài qiǎngjiù tā} — The doctors are fighting to save him.
  \item \emph{Recall:} say the Mandarin for first aid, then the Mandarin for to pass out, then say \emph{qiǎngjiù} again
\end{itemize}

\begin{tcolorbox}[breakable,colback=teal!4,colframe=teal!35!black,title={Before you move on}]
What does qiǎngjiù mean? (\textquotedblleft{}To fight to save a patient; to give emergency treatment\textquotedblright{}.) Say it once more.
\end{tcolorbox}
